\section{数据、资产准入与统计估计}
\subsection{数据范围与资产含义}
资产池由 \etfCount{} 只基金组成，覆盖可转债、国债、信用债、货币市场、中国权益、香港权益、海外权益及商品等风险来源。资产名称和代码以仓库资产定义为准，完整名单列于附录。候选池中的工具不进入本次回测。将候选资产与正式资产分开，可以避免在看到某一市场阶段的表现后临时增补有利工具。

缓存覆盖 2007 年 1 月 18 日至 2026 年 8 月 31 日，正式评价使用 \evalStartDate{} 至 \evalEndDate{}。缓存起点仅表示仓库最早可用观察，不意味着所有资产从该日开始可交易。主模型在评价开始前利用可用历史形成预热，正式比较则统一截取评价区间。

收益来自复权价格序列，保留极端观察。本文不按全样本分布对异常涨跌进行事后删除，因为这会同时改变模型见到的历史和实际承担的损失。是否存在错误报价应依据原始数据证据单独检查，不能把对结果不利的观察自动视作错误。现有价格资料尚不能代替逐笔成交、盘口和基金申赎记录。

\subsection{逐期资产准入}
每个调仓日先截取严格早于该日的最近 \lookbackDays{} 行收益。资产至少具有 60 个有效收益观察且窗口方差为正，才进入合格集合。该门槛用于确保基本估计条件，不属于预测信号。某只基金当期收益较差，并不会因此失去下一期准入资格。

形式上，可以将合格集合写为
\[
\mathcal A_t=\{i:N_{i,t}\geq60,\ \widehat{\operatorname{Var}}(r_i)>0\},
\]
其中 $N_{i,t}$ 只统计调仓日前窗口中的有效观察。上市前的缺失值保留为缺失，不能用未来价格向前填充。对晚上市资产，第一次形成目标权重必须晚于其满足数据条件的日期。这一关系由测试和逐期准入记录检查。

\subsection{共同观察与缺失处理}
均值和协方差使用不同的缺失处理机制。收益估计按各资产有效观察更新，协方差要求跨资产共同完整行。协方差模块先检查每个资产至少有基本观察，再删除含缺失值的整行；若共同观察少于两个，估计失败。该行为不能简化为每期都有 \lookbackDays{} 个完整共同样本。

假设一只新资产已有 60 个有效收益，其他资产已有 \lookbackDays{} 个观察，则该新资产可以准入，但共同完整矩阵可能明显短于 \lookbackDays{} 行。协方差的有效样本量需要从诊断读取，不能直接用窗口上限替代。收缩强度也应根据实际共同矩阵估计，不能先对全样本估计后再套用到历史各期。

在收益记账阶段，现有程序对缺失的单资产日收益不进行持仓重新归一化。指定路径的独立核验显示，持有资产未出现缺失收益，因此该处理没有在评价路径中产生未计价持仓。这个结论仅适用于已核验样本，若未来出现停牌或缺失报价，仍需区分暂缺价格与真实零收益。

\subsection{20 日半衰期收益估计}
令 $\rho=2^{-1/20}$，按有效观察从新到旧编号 $k=0,1,\ldots$。在有限窗口内，收益估计可写为
\[
\widehat\mu_{i,t}=\annualizationDays\,
\frac{\sum_k\rho^k r_{i,t(k)}}{\sum_k\rho^k}.
\]
这里 $t(k)<t$，且只遍历有效观察。程序采用指数加权平均的归一化形式，并跳过缺失值。因而“20 日”在正常连续交易样本中对应 20 个交易日；若资产存在缺失，衰减的有效观察序列与纯日历距离并不完全相同。

半衰期意味着相隔 20 个有效观察的原始权重相差一半。它描述历史信息进入均值的速度，不是持有期约束。组合每周求解一次，但收益估计在每个求解时点使用整个有效窗口。均值年化只是统一单位，不意味着近期收益可以无误差地外推一年。

这一估计会对近期收益变化作出响应。若某项资产近期上涨，收益短缺项可能倾向增加其权重；但参考跟踪与方差项同时存在，最终配置并不等于按近期涨幅排序。反之，负收益估计也不直接产生做空，因为可行域限制为多头预算集合。

\subsection{Ledoit--Wolf 收缩协方差}
对当期共同完整收益矩阵，先构造经验协方差 $S_t$。令
\[
\tau_t=\frac{\operatorname{tr}(S_t)}{n_t},\qquad
\widetilde\Sigma_t=(1-\delta_t)S_t+\delta_t\tau_t I.
\]
实现使用 Ledoit--Wolf 估计器，经验协方差的归一化遵循该库的实现，收缩强度 $\delta_t$ 根据当前矩阵估计。之后乘以 \annualizationDays{} 转换为年化单位，并执行数值半正定修复。本文不把 $\delta_t$ 设为固定的主观比例。

对任意特征向量，收缩会将样本特征值向共同方差水平移动。小特征值不再完全依赖偶然的样本线性关系，大特征值也受到结构化目标约束。目标矩阵为各向同性并不表示最终所有资产方差相同，因为只要收缩强度小于一，样本矩阵仍然参与估计。

\subsection{年化单位与数值修复}
收益均值与协方差均按 \annualizationDays{} 日转换，波动则相应乘以 $\sqrt{\annualizationDays}$。在最终目标中，方差分子和方差尺度使用同一单位，收益短缺分子和收益尺度也一致。若只有一部分输入年化，目标项之间的实际权重会随之改变。

协方差修复包括对称化和尺度相关的特征值下限。修复量必须可追溯，因为参考预算是相对于修复后的矩阵定义的。轻微数值不对称与严重不合法矩阵不是同一问题，程序不能将任意错误矩阵都解释成可接受的噪声。指定结果保存估计方法、回退标记、特征值和求解状态，为这一判断提供依据。

\subsection{数据边界与可推广性}
境内上市海外基金的收益并不完全等同于境外标的的同日收益。时区、汇率及折溢价都可能进入观察结果。本文没有将这些差异单独识别成因子，因此也不据此声称掌握了海外资产的纯本币风险。相同原因下，商品类基金的价格表现不能简单解释为现货价格的机械复制。
